%!TEX TS-program = xelatex
%!TEX encoding = UTF-8 Unicode
% Awesome CV LaTeX Template for Cover Letter
% https://github.com/posquit0/Awesome-CV

%-------------------------------------------------------------------------------
% CONFIGURATIONS
%-------------------------------------------------------------------------------
\documentclass[11pt, a4paper]{awesome-cv}

\geometry{left=1.4cm, top=.8cm, right=1.4cm, bottom=1.8cm, footskip=.5cm}
\fontdir[fonts/]
% --- Use Roboto everywhere (body + headers) ---
\setmainfont{Roboto}
\setsansfont{Roboto}
\renewcommand{\familydefault}{\sfdefault} % optional: enforce sans as default

% --- ALL CAPS for CV section titles ---
\renewcommand*{\sectionstyleface}[1]{%
  {\fontsize{13pt}{1em}\bodyfont\bfseries\MakeUppercase{#1}}%
}
% Color theme
\definecolor{DarkBlue}{HTML}{140097}
\colorlet{awesome}{DarkBlue}
%\colorlet{awesome}{awesome-red}
\setbool{acvSectionColorHighlight}{true}
\renewcommand{\acvHeaderSocialSep}{\quad\textbar\quad}

%-------------------------------------------------------------------------------
% PERSONAL INFORMATION
%-------------------------------------------------------------------------------
% Photo is optional; comment out if you don't want it
% \photo[circle,noedge,left]{profile}

\name{Ivan}{Ovinnikov}
\position{Research Engineer{\enskip\cdotp\enskip}Robust ML (RL, Post-Training, Simulation)}
\address{Zürich, Switzerland}

% Fill these with your actual details (or comment out)
\mobile{(+41) 76 470 89 34}
\email{ivan.ovinnikov@gmail.com}
\homepage{ovinnikov.com}
\github{iov9000}        % e.g. iovinnikov
\linkedin{linkedin-ivan-ovinnikov-0b227593}    % e.g. ivan-ovinnikov
% \googlescholar{scholar-id}{Google Scholar}

%\quote{``Build systems that remain reliable in the long tail."}

%-------------------------------------------------------------------------------
% LETTER INFORMATION
%-------------------------------------------------------------------------------
\recipient
  {Hiring Manager}
  {Google Deepmind}
\letterdate{\today}
\lettertitle{Application for Research Engineer (Visual Design Intelligence)}
\letteropening{Dear Hiring Manager,}
\letterclosing{Sincerely,}
\letterenclosure[Attached]{Curriculum Vitae}

%-------------------------------------------------------------------------------
\begin{document}

\makecvheader[R]

\makecvfooter
  {\today}
  {Ivan Ovinnikov~~~·~~~Cover Letter}
  {}
\makelettertitle

%-------------------------------------------------------------------------------
% LETTER CONTENT
%-------------------------------------------------------------------------------
\begin{cvletter}

I am applying for the Research Engineer role on the Gemini Visual Design Intelligence team. Building a ``visual designer'' requires models that generate editable structure under layout constraints and behave predictably under iterative edits such as reflow, resizing, and element replacement. This is as much a reliability problem as a generative modeling problem: outputs must preserve structure, satisfy constraints, and remain editable under transformation.

My work has focused on closing this capability-to-reliability gap, from safety-critical learning systems at ANYbotics to structured generative modeling at Disney Research.

At ANYbotics, I develop parts of the reinforcement learning stack for quadruped locomotion, including large-scale GPU training with IsaacLab. My focus is reducing the sim-to-real gap through structured curriculum learning: designing environment distributions and domain randomization schemes that expand coverage of contact regimes, terrain variation, and rare failure modes. In parallel, I build and maintain evaluation suites that measure policy consistency across diverse test scenarios and track regressions under controlled distribution shifts rather than optimizing only average-case returns.

During my PhD at ETH Zürich, I built inverse reinforcement learning frameworks for surgical simulation for assistance and skill evaluation. A central challenge was learning expert behavior from sparse demonstrations and inferring intent well enough for a learned policy to act as a training assistant (guidance, feedback, skill assessment). Methodologically, I worked with causal-invariance regularization and optimal-transport--style distribution matching to stabilize learning under shifts in users, tasks, and simulator conditions.

At Disney Research, I worked on generative sequence modeling with structured latent variables, focusing on controllability and structural consistency, which shaped how I think about making generative models predictable and composable.

I am comfortable owning modeling improvements end-to-end, from hypothesis formulation and experiment design to implementation, documentation, and reproducible evaluation. I am also used to working across research, engineering, and product stakeholders, and have experience mentoring and supporting others through teaching and supervision.

Thank you for considering my application. I would be glad to discuss the role further.


\end{cvletter}

\makeletterclosing

\end{document}
